% Extracto editor-ready del contenido exclusivo de masas/08_composicion.tex.
% Las líneas 1--52 son byte-exactas a 75_rev2_ontologia_particula.tex:231--282
% y las líneas 53--396 son byte-exactas a 74_rev2_realizaciones_i.tex:641--1052
% en dos tramos contiguos. Sus propietarios íntegros permanecen preservados;
% esta residencia publica una sola vez el contenido exacto y conserva desde
% la línea 397 todo el desarrollo exclusivo del propietario de composición.

\subsection{Objeto compuesto como sección persistente de un producto cocíclico de rutas}
Una partícula compuesta es una sección del producto fibrado de biografías constituyentes con una frontera de ligadura. Los registros genealógicos se componen antes de aplicar el extractor de firma. Si \(\Gamma\) y \(\Lambda\) se enlazan mediante la frontera \(B\) y la acción \(g\),

\[
(n,e)_\Gamma\star_{(g,B)}(n,e)_\Lambda
=
(n_\Gamma+n_\Lambda-n_B,
e_\Gamma+ge_\Lambda-e_B).
\]
La firma compuesta es

\[
\operatorname{Sig}_{\rm comp}
=L_{\rm tick}\circ\mathfrak C_{12}.
\]
No coincide en general con la suma de las firmas individuales, porque el registro genealógico de enlace porta energía, acarreo y memoria. Esta es la residencia de la energía de ligadura.

\Needspace{7\baselineskip}
\subsection{Resonancia como sección transitoria y anchura como tasa de desintegración}
Una resonancia es una sección metastable cuyo resolvente posee un polo de energía complejo. Si \(z_X=M_Xc_{\rm int}^{2}-i\Gamma_X/2\), la masa es \(M_X=c_{\rm int}^{-2}\operatorname{Re}z_X\); la anchura mide su vida media:

\[
\Gamma_X=\frac{\hbar_{\rm int}}{\tau_X^{\rm life}}.
\]
La anchura no es una sexta coordenada de \(\sigma\). Procede del operador de supervivencia, de los canales de decaimiento y de la continuación analítica del resolvente. El inventario de 384 anchuras exige por ello un lector propio, separado del inventario de 471 masas.

\Needspace{7\baselineskip}
\subsection{Sección virtual como extensión intermedia no asintótica del atlas}
Una sección virtual es una extensión finita que participa en una composición o propagador pero no define una rama global estable en el codominio físico. Si \(L(s)\) es su horizonte,

\[
L(s)<\infty,
\qquad
\operatorname{Ext}_{L(s)}(s)\ne\varnothing,
\qquad
\operatorname{Ext}_{L(s)+1}(s)=\varnothing.
\]
La virtualidad es una obstrucción terminal, no una etiqueta documental ni una masa imaginaria. Una resonancia puede ser inestable sin ser virtual en este sentido; una línea interna puede ser virtual sin constituir una nueva especie.

\Needspace{7\baselineskip}
\subsection{Sección de Majorana como punto fijo de la conjugación partícula–antipartícula}
Cuatro construcciones deben mantenerse separadas. La involución de palabras \(C_M\) define la antisección; un espinor de Majorana satisface una condición de realidad; un modo cero Majorana es una excitación localizada de energía cero en un Hamiltoniano topológico; el sector de Ising posee reglas de fusión

\[
\sigma\otimes\sigma=1\oplus\psi.
\]
Una partícula física es Majorana cuando su representación fermiónica admite una autorconjugación compatible con el Hamiltoniano, la localización y el producto interno. La neutralidad eléctrica y la igualdad de una palabra bajo \(C_M\) no bastan. En particular, decidir la naturaleza Majorana de los neutrinos requiere un entrelazador entre el operador neutro, la base de masas y la estructura real fermiónica.

\Needspace{7\baselineskip}
\subsection{Sección anyónica y representación no trivial del grupo de trenzas}
Un anyón es una sección bidimensional cuya representación del grupo de trenzas no se reduce a la permutación bosónica o fermiónica. En el sector Fibonacci,

\[
\tau\otimes\tau=1\oplus\tau,
\qquad
\operatorname{FPdim}(\tau)=\varphi.
\]
En el sector abeliano \(\mathbb Z/6\mathbb Z\), el cociente de devanado posee seis caracteres unidimensionales. Para obtener una partícula física se necesitan matrices \(F\) y \(R\), pentágono, hexágono, Hamiltoniano con brecha y operadores de trenzado. El anillo de fusión exacto determina el codominio topológico, pero no fija por sí solo una masa universal.

\Needspace{7\baselineskip}
\subsection{Clases taquiónica y fantasma como obstrucciones de positividad y estabilidad}
Una partícula persistente posee evolución unitaria, energía acotada inferiormente y masa espectral real en la recta física. Un taquión exigiría un polo con \(m^2<0\) que permaneciera como estado asintótico estable. En la arquitectura HMT--MD, una dirección que abandona el cono de realizaciones positivas o bien retorna a la frontera nula, o bien pierde supervivencia, o bien se manifiesta como inestabilidad de un fondo que debe condensar. No define una sección persistente superlumínica. La imposibilidad se formula como incompatibilidad entre persistencia global, positividad espectral y \(m^2<0\), no como prohibición de parámetros taquiónicos transitorios en una expansión perturbativa.

Un fantasma es un modo de norma negativa. La positividad del producto interno y la conservación unitaria excluyen su realización física. Los fantasmas de Faddeev--Popov pertenecen al complejo de calibre y no al espacio físico de estados; esta distinción debe conservarse.

\Needspace{7\baselineskip}
\section{Composición, resonancias y anchuras}
\Needspace{7\baselineskip}
\subsection{Composición no aditiva de masas mediante la ley cocíclica del enlace}
La Ley General de Masas actúa sobre biografías genealógicas. Cuando varios constituyentes forman un estado ligado, el objeto de entrada no es la lista de sus masas aisladas, sino un registro compuesto que conserva el orden de acoplamiento, la frontera común, el acarreo, la orientación y la memoria. Sean

\[
\Gamma_1,\ldots,\Gamma_n
\]
rutas constituyentes, \(\mathcal T\) un árbol temporal de composición y \(\partial\) la frontera de enlace. El operador dodecafásico de composición se escribe

\[
\mathfrak C_{12}^{(n)}:
(\Gamma_1,\ldots,\Gamma_n;\mathcal T,\partial)
\longmapsto \mathcal L_{\rm comp}.
\]
Sólo después se aplican el lector del registro, el extractor de firma, el carácter y las torres:

\[
\mathcal L_{\rm comp}
\longmapsto \sigma_{\rm comp}
\longmapsto \chi_{\rm comp}
\longmapsto \widehat M_{\rm comp}.
\]
Esta precedencia es esencial. Si se evaluasen primero las masas aisladas y se sumasen después, desaparecerían precisamente los datos que distinguen un estado ligado de un umbral de constituyentes libres. La energía de ligadura no es un término agregado al final: es la imagen dimensional de una diferencia de registro producida por el enlace.

\Needspace{7\baselineskip}
\subsection{Ley cocíclica del enlace}
Para dos registros \(L_1,L_2\), la composición tiene la forma

\[
L_1\star L_2
=L_1+L_2+\mathfrak b(L_1,L_2),
\]
donde \(\mathfrak b\) registra frontera, acarreo y memoria compartidos. La independencia respecto de la parentización exige la identidad cocíclica

\[
\mathfrak b(L_1,L_2)
+\mathfrak b(L_1\star L_2,L_3)
=
\mathfrak b(L_2,L_3)
+\mathfrak b(L_1,L_2\star L_3).
\]
Si esta igualdad falla, las dos maneras de construir un mismo árbol triple producen registros genealógicos distintos y no existe una especie compuesta bien definida. Si se cumple, el cociclo no desaparece: garantiza que su contribución es una propiedad del estado y no un artefacto de la notación.

Una presentación más explícita conserva además la acción de borde. Sea \(e\in\mathbb Z^{12}\) el vector ordenado de acontecimientos y sea \(n\) el contador de acción o de acontecimientos de un registro. Si \((n,e)_\Gamma\) y \((n,e)_\Lambda\) son registros, si \(B\) es el subregistro común con datos \((n_B,e_B)\), y si \(g\in\operatorname{Aut}(\mathbb Z^{12})\) transporta la orientación del segundo registro a la del primero preservando el lector,

\[
(n,e)_\Gamma\star_{(g,B)}(n,e)_\Lambda
=
(n_\Gamma+n_\Lambda-n_B,
e_\Gamma+g e_\Lambda-e_B).
\]
El término \(B\) evita contar dos veces la frontera identificada; la acción \(g\) transporta el segundo registro a la orientación del primero. La firma compuesta es, por tanto,

\[
\operatorname{Sig}_{\rm comp}
=L_{\rm tick}\!\left(
\mathfrak C_{12}^{(n)}
(\Gamma_1,\ldots,\Gamma_n;\mathcal T,\partial)
\right),
\]
y no la suma automática de firmas elementales.

\Needspace{7\baselineskip}
\subsection{Clausura cocíclica de la composición mesónica quark–antiquark}
Una expresión mesónica orientada tiene la forma

\[
M(q,q')=q\otimes(q')^\vee.
\]
El dominio bruto contiene \(6\times6\times3\times4=432\) expresiones al considerar sabor, color compatible y orientación. Esta cardinalidad no es el número de mesones físicos. Antes de individualizar una especie se aplican:

\begin{enumerate}
\item el proyector de singlete de color;
\item el proyector de carga e isospín;
\item la representación de espín y paridad;
\item el árbol temporal de enlace;
\item la condición radial u orbital;
\item el canal de preparación y decaimiento cuando el estado es resonante.
\end{enumerate}
El descriptor físico completo es

\[
\mathcal I_M=
(q\bar q',J^{PC},I,I_3,S,C,B,\mathcal T,\partial,
\Gamma_M,\mathcal L_M,\sigma_M,\eta_M).
\]
Dos estados con los mismos sabores pueden diferir en \(J^{PC}\), excitación radial, momento orbital, hoja o canal, y por ello poseer polos distintos. La masa no se deriva del mero nombre de los constituyentes.

Las evaluaciones publicadas para \(\pi^0,\pi^\pm,K^\pm,K^0\) son exactas una vez fijados los registros históricos declarados. La afirmación más fuerte ---la generación prospectiva de cada registro nominal desde su árbol de enlace--- requiere que \(\mathfrak C_{12}\) reconstruya esos registros sin consultar el valor que será comparado. Ambas capas se conservan juntas: el cálculo condicionado no se borra, y tampoco se presenta como selector de identidad.

\Needspace{7\baselineskip}
\subsection{Clausura trilineal de la composición bariónica y neutralización del carácter de color}
Una expresión bariónica orientada usa tres constituyentes y tres colores. El dominio bruto del censo adoptado contiene

\[
\lvert\mathcal B_{\rm RGB}^{\rm raw}\rvert=1728
\]
expresiones. La proyección física debe imponer de forma simultánea:

\[
P_{\rm singlete\ de\ color},
\qquad
P_{J^P},
\qquad
P_{I,I_3,S,C,B},
\qquad
P_{\rm Pauli}.
\]
El árbol triple conserva qué par se acopló primero y cómo atravesó la frontera el tercer constituyente. La identidad cocíclica garantiza que las parentizaciones equivalentes representan el mismo estado cuando el entrelazador físico las identifica; los canales inequivalentes permanecen como estados distintos.

Las firmas históricas

\[
\sigma_p=(59,0,9,1,0),
\qquad
\sigma_n=(59,1,-34,0,1)
\]
producen las evaluaciones publicadas de protón y neutrón al aplicar el carácter completo. Son controles exactos del lector para esos registros. La explicación física de la diferencia neutrón--protón exige además derivar, antes del contraste, el signo del isospín, la orientación de vacancias, el borde de enlace y la contribución electromagnética en una carta común.

\Needspace{7\baselineskip}
\subsection{Realizaciones compuestas exóticas y condiciones de estabilidad del enlace}
La misma gramática admite más constituyentes sin modificar la ley:

\[
\mathfrak C_{12}^{(n)}:
\prod_{j=1}^{n}\Gamma_j\times\mathcal T_n\times\partial_n
\longrightarrow\mathcal L_{\rm comp}.
\]
Un tetraquark corresponde a \(qq\bar q\bar q\); un pentaquark, a \(qqqq\bar q\); un glueball, a rutas de la representación adjunta de color proyectadas al singlete; un híbrido, a rutas fermiónicas y de calibre reunidas en un mismo árbol. La lista de constituyentes no determina por sí sola la especie. Una clasificación física completa debe incluir color, espín, paridad, conjugación, isospín, topología del árbol y polo del operador acoplado.

Esta formulación discrimina también dos descripciones rivales de un mismo candidato ---molécula hadrónica y estado compacto, por ejemplo---: son dos árboles o dos proyectores sobre un mismo sector de constituyentes. Se vuelven equivalentes únicamente si existe un unitario que entrelace sus operadores, sus canales y sus polos.

\Needspace{7\baselineskip}
\subsection{Estado estable, estado ligado y resonancia}
Conviene separar tres nociones.

Un estado estable es un autovector normalizable del Hamiltoniano físico con energía real y sin canal de decaimiento abierto. Un estado ligado es un autovector por debajo del umbral pertinente; puede ser estable o decaer por otro canal. Una resonancia es un polo de la continuación del resolvente o de la matriz de dispersión. No es un autovector normalizable del Hamiltoniano autoadjunto sobre el eje real.

Para un polo de energía

\[
z_X=E_X-\frac{i}{2}\Gamma_X
=M_Xc_{\rm int}^{2}-\frac{i}{2}\Gamma_X,
\qquad \Gamma_X>0,
\]
la energía de polo es \(E_X=\operatorname{Re}z_X\), la masa de polo es \(M_X=c_{\rm int}^{-2}\operatorname{Re}z_X\) y la anchura energética es \(\Gamma_X=-2\operatorname{Im}z_X\). La Ley General de Masas suministra la rama real condicionada por la identidad y el registro; la continuación abierta de canales determina el desplazamiento complejo y la anchura. En una carta de unidades naturales \(c_{\rm int}=1\), las coordenadas numéricas de energía y masa coinciden, pero sus tipos no se identifican antes de elegir esa carta.

\Needspace{7\baselineskip}
\subsection{Corrección del signo en la lectura discreta del polo}
Sea \(t_0>0\) un paso temporal y supóngase que el modo de un paso tiene autovalor

\[
\lambda=r e^{-i\theta},
\qquad 0<r<1,
\]
con la convención

\[
\lambda=\exp\!\left(-\frac{i}{\hbar_{\rm int}}z t_0\right).
\]
Entonces

\[
z=\frac{\hbar_{\rm int}}{t_0}
\bigl(\theta+i\log r\bigr)
=E-\frac{i}{2}\Gamma
=Mc_{\rm int}^{2}-\frac{i}{2}\Gamma,
\]
de modo que

\[
E=\frac{\hbar_{\rm int}}{t_0}\theta,
\qquad
M=\frac{\hbar_{\rm int}}{t_0c_{\rm int}^{2}}\theta,
\qquad
\Gamma=-\frac{2\hbar_{\rm int}}{t_0}\log r>0.
\]
El signo positivo delante de \(i\log r\) es obligado: como \(\log r<0\), la parte imaginaria de \(z\) es negativa. La forma \(\theta-i\log r\) colocaría el polo en el semiplano superior y describiría crecimiento, no decaimiento.

La probabilidad de supervivencia tras \(N\) pasos es

\[
|\lambda|^{2N}=r^{2N}
=\exp\!\left(-\frac{\Gamma}{\hbar_{\rm int}}Nt_0\right).
\]
Por ello el tiempo de vida asociado es

\[
\tau_X^{\rm life}=\frac{\hbar_{\rm int}}{\Gamma_X}.
\]
Esta identidad demuestra además por qué la anchura no debe introducirse como sexta coordenada de la firma \(\mathbb Z^5\): nace del módulo del modo abierto, no del carácter positivo que organiza la masa.

\Needspace{7\baselineskip}
\subsection{Operador efectivo de canales}
Sea \(P\) el subespacio preparado y \(Q=I-P\) el conjunto de canales eliminados. El complemento se incorpora mediante el operador de Feshbach

\[
H_{\rm eff}(z)
=PHP+PHQ(z-QHQ)^{-1}QHP.
\]
Los polos satisfacen

\[
\det\bigl(z-H_{\rm eff}(z)\bigr)=0.
\]
El primer término conserva el operador interno preparado. El segundo es la autoenergía de los canales: desplaza la parte real y produce una parte imaginaria cuando se abre el continuo. La masa de carácter y la anchura de supervivencia confluyen en el mismo polo, pero proceden de mapas distintos.

Una formulación equivalente usa un proyector de supervivencia distinto del complemento de Feshbach. Sea

\[
S_X=S_X^*=S_X^2
\]
el proyector sobre el subespacio superviviente del estado preparado. El operador comprimido de un paso es

\[
T_X=S_XUS_X.
\]
Supóngase que \(T_X\) posee un autovalor dominante simple \(\lambda_X\), que el estado preparado tiene solapamiento no nulo con su subespacio espectral y que el resto del espectro está separado en módulo. Si \(r_X=|\lambda_X|=\rho(T_X)<1\), entonces

\[
\Gamma_X=-\frac{2\hbar_{\rm int}}{t_0}\log r_X.
\]
Si el operador se normaliza desde el principio sobre probabilidades en vez de amplitudes, desaparece el factor dos. La publicación debe declarar cuál de las dos convenciones usa; mezclarlas produce una anchura errónea por un factor dos.

\Needspace{7\baselineskip}
\subsection{Conjugación de partícula y antipartícula}
Sea \(C_M\) la involución que invierte orientación y hoja, y sea \(U_C\) su realización unitaria sobre el espacio físico. La igualdad de masas de polo y anchuras requiere covariancia tanto del operador como de los canales:

\[
U_CH_{\rm eff}^{X}(z)U_C^{-1}
=H_{\rm eff}^{\bar X}(z).
\]
Entonces los conjuntos de polos retardados coinciden después de transportar el sector y la prescripción causal, y

\[
M_{\bar X}=M_X,
\qquad
\Gamma_{\bar X}=\Gamma_X.
\]
La equivalencia unitaria desnuda no convierte el polo retardado en su conjugado complejo; transporta el mismo problema espectral al sector antipartícula. La paridad de la firma másica no basta para demostrar la segunda igualdad: también los canales y sus medidas deben ser covariantes. Si la simetría se representa antiunitariamente, debe transportarse además la prescripción retardada; la conjugación desnuda \(z\mapsto\bar z\) intercambia polos retardados y avanzados y no es, por sí sola, la afirmación física de igualdad de vidas medias.

\Needspace{7\baselineskip}
\subsection{Inventario tipado de realizaciones y observables de contraste}
El corte externo preservado contiene

\[
471\ \text{observables de masa},
\qquad
384\ \text{observables de anchura}.
\]
La partición de masas es

\Needspace{7\baselineskip}
\begingroup
\footnotesize
\setlength{\tabcolsep}{1.9pt}
\begin{longtable}{L{5.04cm}L{5.04cm}L{5.04cm}}
\toprule
\rowcolor{HMTNavy}\HMTtablehead{Dominio físico de llegada} & \HMTtablehead{Masas} & \HMTtablehead{Anchuras} \\
\midrule
\endfirsthead
\multicolumn{3}{l}{\sffamily\scriptsize\color{HMTGray}Continuación de la tabla}\\[-1pt]
\toprule
\rowcolor{HMTNavy}\HMTtablehead{Dominio físico de llegada} & \HMTtablehead{Masas} & \HMTtablehead{Anchuras} \\
\midrule
\endhead
\midrule
\endfoot
\bottomrule
\endlastfoot
partículas elementales, leptónicas y extensiones & 53 & 9 \\
\rowcolor{HMTLight}quarks y extensiones quark & 8 & 1 \\
mesones & 243 & 225 \\
\rowcolor{HMTLight}bariones & 166 & 149 \\
registro gravitatorio de contraste & 1 & 0 \\
\rowcolor{HMTLight}\textbf{Total} & \textbf{471} & \textbf{384} \\
\end{longtable}
\endgroup
Estas filas son observables, no especies elementales. Varias filas pueden referirse a diferentes cargas de un mismo multiplete, estimadores de una misma resonancia, masas Breit--Wigner, masas de polo o estados excitados. La individualización interna precede siempre a la fila externa.

El inventario conserva además una propiedad metodológica fuerte: ninguna de sus 855 magnitudes fue autorizada como selector de ruta. Los 471 registros de masa y los 384 de anchura sólo se abren después de congelar la identidad HMT, el proyector, el operador y la carta de comparación.

\iffalse
\Needspace{7\baselineskip}
\subsection{Malla histórica extendida de pares y resonancias}
Una formulación histórica asoció a la malla angular extendida la expresión

\[
\Theta(n,k,t)
=nA_{\rm rad}-k\frac{C^*_{\rm rad}}6+t\Delta_4.
\]
Con las definiciones vigentes de \(A\), \(C^*\), \(\Delta_4\) y de la carta exponencial, esta expresión aislada no reproduce las cifras archivadas. Se conserva como hipótesis histórica de parametrización, no como su generador ejecutable. Conservamos las doce evaluaciones importadas porque forman parte de la genealogía numérica de la ley. Cinco nombres ---\(\eta,\rho^0,\phi(1020),J/\psi\) y \(\Upsilon(1S)\)--- amplían materialmente las etiquetas visibles respecto de las diecisiete evaluaciones rectoras:

\Needspace{7\baselineskip}
\begingroup
\fontsize{7.2}{8.2}\selectfont
\setlength{\tabcolsep}{1.9pt}
\begin{longtable}{L{1.25cm}L{0.78cm}L{0.78cm}L{0.78cm}L{3.43cm}L{3.90cm}L{4.06cm}}
\toprule
\rowcolor{HMTNavy}\HMTtablehead{Estado} & \HMTtablehead{\(n\)} & \HMTtablehead{\(k\)} & \HMTtablehead{\(t\)} & \HMTtablehead{Referencia posterior (MeV/\allowbreak{}c\({}^{2}\))} & \HMTtablehead{Evaluación HMT histórica (MeV/\allowbreak{}c\({}^{2}\))} & \HMTtablehead{Diferencia (ppm)} \\
\midrule
\endfirsthead
\multicolumn{7}{l}{\sffamily\scriptsize\color{HMTGray}Continuación de la tabla}\\[-1pt]
\toprule
\rowcolor{HMTNavy}\HMTtablehead{Estado} & \HMTtablehead{\(n\)} & \HMTtablehead{\(k\)} & \HMTtablehead{\(t\)} & \HMTtablehead{Referencia posterior (MeV/\allowbreak{}c\({}^{2}\))} & \HMTtablehead{Evaluación HMT histórica (MeV/\allowbreak{}c\({}^{2}\))} & \HMTtablehead{Diferencia (ppm)} \\
\midrule
\endhead
\midrule
\endfoot
\bottomrule
\endlastfoot
\(\mu\) & 64 & 146 & 1 & 105.658374 & 105.656529 & -17.461938\allowbreak{}2274 \\
\rowcolor{HMTLight}\(\tau\) & 64 & 349 & -1 & 1776.86 & 1776.835 & -14.069763\allowbreak{}5154 \\
\(\pi^\pm\) & 37 & 125 & 0 & 139.57039 & 139.570582 & 1.375649\allowbreak{}94982 \\
\rowcolor{HMTLight}\(K^\pm\) & 52 & 177 & 1 & 493.677 & 493.679 & 4.051231\allowbreak{}87833 \\
\(K^0\) & 53 & 178 & 1 & 497.611 & 497.616 & 10.048009\allowbreak{}3889 \\
\rowcolor{HMTLight}\(p\) & 55 & 195 & 1 & 938.272081 & 938.27244 & 0.382618\allowbreak{}226919 \\
\(n\) & 55 & 197 & 0 & 939.565413 & 939.538 & -29.176254\allowbreak{}9161 \\
\rowcolor{HMTLight}\(\eta\) & 46 & 158 & 0 & 547.862 & 547.836 & -47.457206\allowbreak{}3768 \\
\(\rho^0\) & 50 & 176 & 0 & 775.26 & 775.267 & 9.029228\allowbreak{}90385 \\
\rowcolor{HMTLight}\(\phi(1020)\) & 54 & 190 & 0 & 1019.461 & 1019.475 & 13.732747\allowbreak{}0104 \\
\(J/\allowbreak{}\psi\) & 64 & 237 & 0 & 3096.916 & 3096.979 & 20.342818\allowbreak{}4684 \\
\rowcolor{HMTLight}\(\Upsilon(1S)\) & 89 & 334 & 0 & 9460.3 & 9460.34 & 4.228195\allowbreak{}72318 \\
\end{longtable}
\endgroup
Las doce cifras archivadas son evaluaciones históricas importadas; no se atribuyen a \(\Theta\) hasta localizar la carta que las reproduzca. Además, la malla no construye el morfismo

\[
\text{estado compuesto}
\longrightarrow\text{ruta y registro genealógico}
\longrightarrow(n,k,t).
\]
Además conserva la referencia externa en la misma tabla histórica. Por ello su fuerza es evidencia numérica condicionada, no selección física anterior al contraste. Publicarla amplía el censo visible sin convertir cinco resonancias en cinco predicciones anteriores al contraste ni desplazar las evaluaciones E3 vigentes.
\fi

\Needspace{7\baselineskip}
\subsection{Normalización de identidades nominales en registros comparables}
El empalme prospectivo completo tiene la forma

\[
\begin{aligned}
\text{identidad física}
&\longrightarrow
\text{sector y representación}
\longrightarrow
\text{proyector de fibra o árbol compuesto}\\
&\longrightarrow
\text{ruta enriquecida}
\longrightarrow
\text{registro genealógico}
\longrightarrow
\text{firma y torres}\\
&\longrightarrow
\text{operador preparado}
\longrightarrow
\text{polo o autovalor}
\longrightarrow
\text{carta dimensional}\\
&\longrightarrow
\text{comparación externa}.
\end{aligned}
\]
Para un estado estable, la salida operatoria es un autovalor real. Para una resonancia, es un polo complejo. Para un quark, el mapa incluye esquema y escala. Para un límite experimental, incluye el nivel de confianza. Para un observable correlacionado, incluye la matriz de covarianza. Estas diferencias son estructurales: determinan qué magnitudes pueden compararse.

\Needspace{7\baselineskip}
\subsection{Condiciones de clausura individual de las realizaciones materiales}
Una fila del inventario queda individualizada desde HMT--MD cuando se han construido, sin consultar su magnitud externa:

\begin{enumerate}
\item el sector y la representación;
\item la ruta, multisección o árbol de composición;
\item el registro genealógico y la firma;
\item el operador de masa y, si procede, el operador de canales;
\item el proyector que selecciona el estado nombrado;
\item la convención de masa, anchura, esquema y escala;
\item la carta dimensional;
\item el falsador interno.
\end{enumerate}
La ley universal ya proporciona la gramática común. El trabajo por estado no consiste en inventar una fórmula diferente, sino en construir el entrelazador que transporta una identidad física hasta el operador apropiado.

\Needspace{7\baselineskip}
\subsection{obstrucciones y condiciones de no degeneración específicos}
La composición queda refutada si dos parentizaciones físicamente equivalentes dan firmas o espectros distintos. El selector nominal queda refutado si cambia al ocultar la columna externa de masa. El lector de anchura queda refutado si produce \(\Gamma<0\) para una contracción, confunde amplitud con probabilidad o asigna anchura a una sección estable sin canal abierto. La carta de resonancia queda refutada si no declara masa de polo o Breit--Wigner. La comparación quark queda indeterminada si omite esquema o escala. La igualdad partícula--antipartícula de anchuras queda refutada si el operador de canales no es covariante bajo conjugación.

\Needspace{7\baselineskip}
\subsection{Transporte de las clases materiales al corredor excepcional mediante rutas, firmas y representaciones}
La gramática de composición, los dominios mesónico y bariónico, el operador de Feshbach, el censo 471/384 y el principio de comparación se conservan en el desarrollo integral antecedente, secciones 16 y 18. Los registros individualizados se presentan íntegramente en los inventarios exhaustivos de masas y anchuras del apéndice. La corrección del signo del polo se obtiene directamente de \(\lambda=\exp(-izt_0/\hbar_{\rm int})\) y fortalece la formulación sin alterar ninguno de esos registros.

\Needspace{7\baselineskip}
\section{Sectores topológicos, extensiones hipotéticas y gravedad}
\Needspace{7\baselineskip}
\subsection{Separación tipada entre sectores topológicos, extensiones hipotéticas y realización gravitatoria}
La arquitectura distingue con precisión:

\begin{enumerate}
\item una identidad algebraica interna;
\item un codominio de fusión o trenzado;
\item una partícula física con Hamiltoniano, brecha, localización y canal de
medida.

\end{enumerate}
Una regla de fusión exacta no es todavía una masa. Un modo de borde no es todavía una especie. Una involución de rutas no es todavía una condición de Majorana. El paso físico exige un funtor que preserve producto, involución, orientación, energía y composición de caminos.

\Needspace{7\baselineskip}
\subsection{Sector de Fibonacci: anillo basado, fusión y representación de trenzas}
La incidencia areal--volumétrica está gobernada por

\[
F_{\rm av}=
\begin{pmatrix}
0&1\\[2pt]1&1
\end{pmatrix}.
\]
Su polinomio característico es

\[
x^2-x-1,
\]
y su autovalor dominante es \(\varphi\). En la base \(\{\mathbf1,\tau\}\), la multiplicación por \(\tau\) se representa por \(F_{\rm av}\); por tanto

\[
\tau\otimes\tau\cong\mathbf1\oplus\tau,
\qquad
\operatorname{FPdim}(\tau)=\varphi.
\]
Esta es una realización exacta del anillo basado de Fibonacci. Para obtener una categoría anyónica física se deben construir además el asociador \(F\), el trenzado \(R\), las ecuaciones de pentágono y hexágono, la quiralidad, un Hamiltoniano local con brecha y el mapa de preparación y lectura. La salida de masa, cuando existe, es una brecha de energía convertida por una velocidad efectiva:

\[
m_{\rm brecha}=\frac{\Delta E}{c_{\rm eff}^{2}}.
\]
No existe una masa universal del anyón Fibonacci deducible únicamente de \(\varphi\).

\Needspace{7\baselineskip}
\subsection{Separación entre el sector topológico y el centro apuntado del atlas}
El centro trenzado del sector apuntado basado en \(\mathbb Z/9\mathbb Z\) posee 81 objetos simples invertibles, cada uno de dimensión cuántica uno. Un simple Fibonacci exigiría dimensión \(\varphi\), y un simple Ising exigiría \(\sqrt2\). Por tanto ninguno de los dos aparece como objeto simple interno de ese centro apuntado.

Este control no elimina las realizaciones topológicas. Determina su lugar: deben obtenerse mediante extensión, condensación, cociente, defecto o funtor monoidal hacia un codominio no apuntado. Presentar el anillo Fibonacci como si ya fuese toda la categoría borraría precisamente el morfismo físico que se debe construir.

\Needspace{7\baselineskip}
\subsection{Sector Ising y modos de Majorana}
La categoría de Ising tiene objetos simples

\[
\mathbf1,\quad\psi,\quad\sigma
\]
y reglas

\[
\psi\otimes\psi=\mathbf1,
\qquad
\psi\otimes\sigma=\sigma,
\qquad
\sigma\otimes\sigma=\mathbf1\oplus\psi,
\qquad
d_\sigma=\sqrt2.
\]
Cuatro objetos distintos suelen recibir nombres próximos y deben conservarse separados:

\begin{itemize}
\item la involución de rutas \(C_M=-\operatorname{rev}\);
\item la paridad de la completación CAR;
\item un espinor relativista autorconjugado;
\item un modo cero topológico de Majorana en una fase Ising.
\end{itemize}
Un fermión relativista es de Majorana cuando existe una conjugación real \(\mathcal C\) compatible con su representación y dinámica,

\[
\psi=\mathcal C\bar\psi^{T},
\]
y el término de masa respeta esa condición. Un modo cero de Majorana se describe mediante un operador autoadjunto \(\gamma=\gamma^\dagger\) que satisface

\[
\{\gamma_i,\gamma_j\}=2\delta_{ij}
\]
y permanece separado del continuo por una brecha. La igualdad de una palabra con su imagen por \(C_M\) no demuestra ninguna de estas dos realizaciones.

Para los neutrinos, la alternativa Dirac/Majorana sólo se decide después de construir el operador neutro, su base de masas, el entrelazador de mezcla y la estructura real fermiónica. La neutralidad eléctrica es necesaria para una autorconjugación no trivial, pero no suficiente.

\Needspace{7\baselineskip}
\subsection{Trenzado abeliano de orden \(6\)}
Sea

\[
\operatorname{wind}_6:B_n\longrightarrow\mathbb Z/6\mathbb Z
\]
el lector de devanado y

\[
\chi_k(b)=\exp\!\left(
\frac{2\pi i k}{6}\operatorname{wind}_6(b)\right).
\]
Las realizaciones con fase \(q\) descienden a meras permutaciones sólo cuando

\[
q^2=1.
\]
Los cuatro caracteres restantes distinguen clases algebraicas del lector de devanado. El resultado exacto es la clasificación de los seis caracteres unidimensionales del cociente de devanado; no es todavía una clasificación de seis fases materiales. Para afirmar que conservan trenzado físico se requieren operadores de trenza y una representación del grupo de trenzas. Su existencia no implica cuatro nuevas partículas: una especie anyónica exige un Hamiltoniano, operadores de trenza, brecha, estabilidad y un acoplamiento a una multisección física. Su promoción a fases materiales depende de esa realización.

\Needspace{7\baselineskip}
\subsection{Positividad y exclusión del taquión persistente}
La recta de acción posee orientación positiva,

\[
\hbar_{\rm int}>0,
\]
el carácter de masa toma valores positivos,

\[
\chi_{\rm full}(\sigma,\eta)>0,
\]
y el refinamiento bidireccional se implementa por conjugación. Supóngase explícitamente

\[
\widehat M^{(0)}\ge0,
\qquad
\widehat K=\widehat K^*,
\qquad
R_{\rm act}>0.
\]
Entonces

\[
\widehat M^{\beta}
=R_{\rm act}^{\widehat K/2}
\widehat M^{(0)}
R_{\rm act}^{\widehat K/2}\ge0.
\]
En consecuencia,

\[
\operatorname{Spec}(\widehat M^\beta)\subset[0,\infty).
\]
Por tanto el operador de masa al cuadrado

\[
\widehat{M^2}_{\beta}:=(\widehat M^\beta)^2
\]
es positivo y satisface

\[
\operatorname{Spec}(\widehat{M^2}_{\beta})\subset[0,\infty).
\]
Una partícula asintótica persistente que satisfaga las premisas de positividad y unitariedad no puede tener \(m^2<0\). Si la linealización de un fondo produce frecuencia imaginaria, el objeto es una dirección inestable: el fondo debe relajarse o condensar. Si la ruta deja de prolongarse, el objeto es una sección virtual con obstrucción terminal. Ninguno de los dos casos define una partícula estable superlumínica.

El alcance de la proposición está fijado por sus premisas. Un parámetro "taquiónico" en un potencial perturbativo puede señalar una transición de fase; lo excluido es su promoción a especie persistente de norma positiva.

\Needspace{7\baselineskip}
\subsection{Exclusión independiente de fantasmas}
Un fantasma físico sería un modo de norma negativa. Sea \(\mathcal H_{\rm phys}=\ker G/\operatorname{im}G_0\) el cociente de calibre dotado de producto interno positivo, y sea \(\widehat M_{\rm phys}\) autoadjunto sobre él. Entonces la medida espectral de un estado preparado satisface

\[
\mu_{P,a}^{\rm mass}(B)\ge0.
\]
Una dirección negativa indica que falló la positividad del cociente o que el modo no pertenece al espacio físico; no desaparece por una afirmación nominal. Los fantasmas de Faddeev--Popov son variables del complejo de calibre, no partículas asintóticas. Esta exclusión es distinta de la taquiónica: un fantasma viola positividad de norma; una dirección taquiónica viola estabilidad del fondo.

\Needspace{7\baselineskip}
\subsection{Fotón, sesgo unidireccional y puente Proca--Stückelberg}
Los canales armónicos producen el sesgo

\[
\varepsilon=q_-^{90}-q_+^{90}.
\]
Una realización cinemática bidireccional puede escribirse

\[
c_+=\frac{c_{\rm int}}{1+\varepsilon},
\qquad
c_-=\frac{c_{\rm int}}{1-\varepsilon},
\]
para la cual la medida de ida y vuelta conserva exactamente

\[
\frac{2}{c_+^{-1}+c_-^{-1}}=c_{\rm int}.
\]
El sesgo permanece en la lectura unidireccional y cancela en la lectura recíproca. Esta identidad pertenece a la arquitectura de dos orientaciones; no demuestra por sí sola una masa fotónica.

Si el cierre de frontera produce una frecuencia de brecha \(\mu_T\in\mathcal L_T^{-1}\) en el registro,

\[
\omega^2=c_{\rm int}^{2}k^2+\mu_T^2,
\]
un entrelazador dimensional puede transportarlo a una realización Proca--Stückelberg con masa efectiva

\[
m_\gamma=\frac{\hbar_{\rm int}\mu_T}{c_{\rm int}^{2}}.
\]
Para una propagación de frecuencia \(\omega\), la corrección temporal de primer orden satisface

\[
\frac{\Delta T}{T}
\simeq
\frac12
\left(
\frac{m_\gamma c_{\rm int}^{2}}
     {\hbar_{\rm int}\omega}
\right)^2.
\]
La cadena requerida es

\[
\text{ida/retorno TPK}
\longrightarrow
\text{memoria}
\longrightarrow
\mu_T
\longrightarrow
m_\gamma.
\]
No se identifica automáticamente \(\varepsilon\), el cociente de acción \(R_{\rm act}\) y \(\mu_T\): viven en codominios distintos. La carta nula del fotón ordinario permanece exacta mientras no se construya la brecha y su entrelazador físico.

\Needspace{7\baselineskip}
\subsection{Clasificación operatoria de los sectores oscuros por dominio, lector y codominio físico}
Un sector oscuro no se define por una masa arbitraria, sino por el núcleo de sus acoplamientos. Si \(P_X\) es su proyector,

\[
P_{\rm em}P_X=0,
\qquad
P_{\rm color}P_X=0,
\]
mientras puede ocurrir

\[
P_{\rm grav}P_X\ne0
\]
o existir un acoplamiento débil. La ley entrega primero un operador de fibra; el canal determina después si su espectro contiene un modo estable, metaestable o resonante.

Un fotón oscuro requiere una segunda conexión \(A'\), su curvatura \(F'\) y un entrelazador cinético

\[
\mathcal L_{\rm mix}
=-\frac{\varepsilon_{\rm kin}}2F_{\mu\nu}F'^{\mu\nu}.
\]
Si el segundo sector adquiere masa por Stückelberg, Higgs oculto o frontera, su polo es distinto del fotón nulo. La magnitud \(\varepsilon_{\rm kin}\) debe derivarse de una incidencia común antes de comparar observables.

Sean \(P_{\rm weak},P_{\rm em},P_{\rm color}\) los proyectores sobre los canales de representación débil, electromagnético y cromático, y \(P_X\) el soporte de la sección oscura. Un WIMP exige solapamiento débil no nulo y ausencia de soporte electromagnético y cromático:

\[
P_{\rm weak}P_X\ne0,
\qquad
P_{\rm em}P_X=P_{\rm color}P_X=0.
\]
El primer producto expresa incidencia de subespacios; un conmutador no nulo no sería un criterio de acoplamiento, pues dos proyectores pueden conmutar y, aun así, compartir soporte. La intensidad de la interacción requiere después un operador de vértice o un Hamiltoniano de acoplamiento.

Su estabilidad requiere un autovalor unitario del transporte o una paridad protectora. "WIMP" es una condición de representación y persistencia, no una masa predeterminada.

\Needspace{7\baselineskip}
\subsection{Clasificación de la sección neutrínica estéril por fibra invariante, mezcla y carga de calibre nula}
Un neutrino estéril es un autovector del bloque neutro que satisface

\[
P_{\rm weak}n_s=0,
\qquad
P_{\rm mass}n_s=n_s.
\]
No es sinónimo de la cuarta profundidad estructural del soporte neutrínico. Su existencia requiere un núcleo débil no nulo, un autovalor del operador de masa en ese núcleo y un transporte de mezcla con los modos activos. El criterio es operatorio y puede falsarse antes de introducir una escala externa.

\Needspace{7\baselineskip}
\subsection{Clasificación de las secciones axiónica y pseudoescalares por carácter topológico y acoplamiento de borde}
Un candidato axiónico debe ser un modo pseudoescalar con término cinético normalizado,

\[
CP\,a\,CP^{-1}=-a,
\qquad
\mathcal L_{\rm kin}=\frac{Z_a}{2}\,\partial_\mu a\,\partial^\mu a,
\qquad Z_a>0,
\]
acoplado a la densidad topológica de color. Su masa procede de la curvatura del potencial efectivo,

\[
m_a^2=\frac1{Z_a}
\left.
\frac{\partial^2V_{\rm eff}}{\partial a^2}
\right|_{a=a_0}.
\]
El canal \(C^*\mapsto-C^*\) aporta una orientación impar, pero no basta para construir \(V_{\rm eff}\), el término cinético ni el proyector físico. La especie sólo queda definida cuando esos mapas están presentes.

\Needspace{7\baselineskip}
\subsection{Realizaciones monopolares y defectos topológicos: clases de carga y estabilidad}
Sea \(F\) la curvatura de una conexión \(U(1)\) en la normalización integral elegida. Un monopolo magnético se caracteriza por la clase

\[
\frac1{2\pi}\int_{S^2}F=n\ne0.
\]
HMT debe realizar \(n\) como holonomía neta de una familia cerrada de rutas y demostrar que sobrevive al refinamiento. Una vacancia puntual o un acarreo aislado no son un monopolo. Si el defecto es dinámico y estable, su masa es un autovalor del operador linealizado alrededor de esa clase.

\Needspace{7\baselineskip}
\subsection{Clasificación de las extensiones supersimétricas mediante emparejamiento de representaciones y caracteres}
La coexistencia de completaciones fermiónica y bosónica no demuestra supersimetría. En cuatro dimensiones se necesita una familia de cargas impares que satisfaga

\[
Q_\alpha:\mathcal F_-\longleftrightarrow\mathcal F_+,
\qquad
\{Q_\alpha,\bar Q_{\dot\beta}\}
=2\sigma^\mu_{\alpha\dot\beta}P_\mu,
\]
que entrelace los operadores de masa:

\[
Q_\alpha\widehat M_-=\widehat M_+Q_\alpha.
\]
La identidad \(Q^2=H\) se reserva para una reducción hermítica unidimensional explícita; no sustituye al álgebra relativista anterior.

Sólo entonces aparecen pares espectrales. La ruptura debe derivar la separación de masa sin usar el valor del compañero buscado. Selectrones, neutralinos, gluinos u otros nombres permanecen como realizaciones posibles, no como filas predichas por la mera existencia de CAR y CCR.

\Needspace{7\baselineskip}
\subsection{Clasificación de las secciones de dilatón y radión mediante modos escalares de volumen y compactificación}
El dilatón es una fluctuación de la sección de escala; el radión, de una longitud o volumen interno. Sus acciones elementales son

\[
t_0\mapsto e^{\phi_d}t_0,
\qquad
\mathcal V_{\rm int}\mapsto e^{\phi_r}\mathcal V_{\rm int}.
\]
Para constituir partículas necesitan término cinético, potencial, proyector persistente y polo. Antes de ello son grados colectivos de la carta dimensional. Esta distinción evita convertir una libertad de normalización en una nueva especie escalar.

\Needspace{7\baselineskip}
\subsection{Compuestos y energía de enlace}
Un estado compuesto \(Y\) con constituyentes \(X_1,\ldots,X_k\) se construye en el orden

\[
(\Gamma_{X_1},\ldots,\Gamma_{X_k})
\longrightarrow
\mathscr L_{\rm bind}
\longrightarrow
\sigma_Y
\longrightarrow
\eta_Y
\longrightarrow
\widehat M_Y.
\]
El registro de enlace \(\mathscr L_{\rm bind}\) contiene el orden de las rutas, la frontera común, el acarreo transferido, el canal de color y la holonomía de ligadura. No es lícito reemplazarlo por \(\sum_i m_{X_i}\). Mesones y bariones requieren lectores distintos porque sus números de constituyentes, representaciones y simetrías de intercambio son diferentes.

Las seis evaluaciones compuestas publicadas prueban la evaluación del carácter al fijar la firma. El cierre estado por estado exige además que el registro de enlace regenere esa firma. Esta condición se aplica del mismo modo a núcleos, átomos, moléculas y candidatos tetraquark o pentaquark, con sus respectivos Hamiltonianos de frontera.

\Needspace{7\baselineskip}
\subsection{Anchura como 2.º observable}
Las 384 anchuras no se obtienen añadiendo una sexta coordenada a \(\sigma\). Para cada estado inestable se construyen los canales \(\mathcal H_{\rm in}\oplus\mathcal H_{\rm out}\), el operador efectivo \(H_{\rm eff}(z)\) y el polo

\[
z_X=M_Xc_{\rm int}^{2}-\frac{i}{2}\Gamma_X.
\]
La parte real pertenece a la ley de masa; la parte imaginaria, al lector de decaimiento. La relación \(\tau_X=\hbar_{\rm int}/\Gamma_X\) se interpreta después de fijar la carta temporal. Así, una teoría exhaustiva de 471 masas y 384 anchuras contiene dos lectores coordinados y no una sola columna con 855 números heterogéneos.

\Needspace{7\baselineskip}
